\section{Player-Level Models}
\label{sec:player_level_models}

This section extends the P7 framework from the club level to the player level, decomposing team strength into individual contributions.
While club-level models treat teams as atomic entities, the player-level specification allocates team performance to the players who participated in each match.
In this model, each player is assigned a latent log-skill parameter $\alpha_p$, constrained under a sum-to-zero identifiability condition.
This constraint ensures that estimated abilities represent deviations from the league-wide average player, rather than absolute or context-free measures of talent.

To evaluate how temporal aggregation affects estimation, we consider two time periods.
The first one corresponds to the 2025 season, the latest full season of the \textit{Brasileirão}.
The second one covers the 2020--2025 period, analysed in Section~\ref{sec:club_applications}.
For each period, the model estimates individual log-skill parameters for all players who participated in at least one match, along with three global parameters.
The global parameters measure the home effect ($\nu$), the common effect parameter ($\eta$), and the impact of a red card by the introduction of a synthetic player parameter ($\alpha_{\text{synthetic}}$).

Player rankings are reported using posterior medians of $\alpha_p$ for each player.
For descriptive purposes, we also report total minutes played, number of matches, wins, draws, and a set of weighted metrics based on playing time.
In these weighted metrics we use the fraction of minutes played by the player in the match to compute their impact on points, goals, and goals difference.
For example, a player who played 30 minutes of a match receives 33.3\% of the team's points, goals, and goals difference for that match.
Also, to help identify the player, the column ``Team'' reports the team where the player played the most minutes during the analysed period.

\subsection{2025 Season}
\label{subsec:player_2025}

Section~\ref{subsec:season_2025} analysed the 2025 \textit{Brasileirão} from a club-level perspective, showing how Flamengo and Palmeiras contested the title throughout the season, with Flamengo securing the championship in the penultimate round.
The relegation battle extended to the final match-day, with Ceará and Fortaleza joining Sport and Juventude in the bottom four.
Here, we extend this analysis to the player level, decomposing the 20 teams' performance into 730 unique players participating across the 380 matches.
The first point of interest is given by the global parameters, in Table~\ref{tab:player_global_2025}.

\begin{table}[htbp]
    \centering
    \begin{tabular}{lccccc}
        \toprule
        Parameter & Mean & Median & Std & 95\% CI \\
        \midrule
        $\nu$ (home effect) & 0.422 & 0.422 & 0.066 & [0.294, 0.549] \\
        $\eta$ (common effect) & $-0.002$ & $-0.002$ & 0.051 & [$-0.104$, 0.096] \\
        $\alpha_{\text{synthetic}}$ & $-2.51$ & $-2.28$ & 1.69 & $(-\infty, -0.16]$ \\
        \bottomrule
    \end{tabular}
    \caption{
        \textbf{Posterior estimates for global parameters under the P7 model for the 2025 season}.
        Estimates are computed from 30,000 MCMC samples.
    }
    \label{tab:player_global_2025}
\end{table}

The home effect estimate of $\nu = 0.422$ corresponds to a multiplicative factor of $\exp(0.422) \approx 1.52$ on expected goals for the home team.
This value is nearly identical to the club-level estimate reported in Table~\ref{tab:match_params} ($\nu = 0.425$), showing consistency between the two levels of analysis.
The common effect parameter $\eta \approx 0$ indicates a negligible systematic deviation in overall scoring levels beyond what is explained by team strengths.
This value is slightly different from the one estimated at the club-level ($\eta = -0.062$), but both credible intervals include zero.

The new parameter representing the synthetic player, $\alpha_{\text{synthetic}} = -2.28$, quantifies the effect of playing with fewer players due to red cards.
Since we are interested in answering whether a red card has a negative impact on club performance, we calculate a unilateral credible interval.
Given the credible interval, we can affirm that a red card acts as a below-average player.
Compared with other parameters, this estimate has a larger variance because the sum-to-zero constraint defines $\alpha_{\text{synthetic}} = -\sum_{p} \alpha_p$, causing its uncertainty to accumulate across all individual player estimates.

Beyond global parameters, the model estimates individual skill coefficients for each player.
These estimates capture each player's contribution to their team's aggregate strength, as defined by Equation~\ref{eq:team_skill_aggregation}.
This aggregate strength determines the expected goal rate relative to the opponent's strength, after accounting for home effect.
Table~\ref{tab:top20_2025} presents the top 20 players by estimated log-skill.

\begin{table}[htbp]
    \centering
    \tiny
    \begin{tabular}{rlcccrrrrrrrr}
        \toprule
        \# & Player & Team & Median & 95\% CI & Min. & G & W & D & WS & WG & WGD \\
        \midrule
        1 & Agustin Rossi & Flamengo & 0.049 & [-0.15, 0.25] & 3,330 & 37 & 23 & 9 & 78.0 & 75.0 & 51.0 \\
        2 & Léo Pereira & Flamengo & 0.046 & [-0.16, 0.25] & 2,954 & 33 & 22 & 8 & 73.5 & 68.5 & 47.5 \\
        3 & De Arrascaeta & Flamengo & 0.037 & [-0.16, 0.23] & 2,402 & 33 & 22 & 7 & 58.5 & 52.2 & 35.8 \\
        4 & Léo Ortiz & Flamengo & 0.030 & [-0.17, 0.23] & 2,248 & 26 & 16 & 7 & 52.0 & 45.0 & 31.7 \\
        5 & Luiz Araújo & Flamengo & 0.029 & [-0.17, 0.23] & 1,712 & 33 & 22 & 7 & 44.0 & 40.9 & 29.3 \\
        6 & Villalba & Cruzeiro & 0.027 & [-0.17, 0.23] & 2,864 & 32 & 16 & 13 & 60.5 & 47.3 & 27.3 \\
        7 & Weverton & Palmeiras & 0.027 & [-0.17, 0.23] & 2,340 & 26 & 19 & 3 & 60.0 & 49.0 & 28.0 \\
        8 & Jemmes & Mirassol & 0.026 & [-0.17, 0.22] & 3,318 & 37 & 18 & 13 & 66.6 & 62.7 & 25.9 \\
        9 & Lucas Silva & Cruzeiro & 0.025 & [-0.17, 0.22] & 2,752 & 34 & 18 & 12 & 58.5 & 45.9 & 24.6 \\
        10 & Christian & Cruzeiro & 0.025 & [-0.17, 0.22] & 2,387 & 30 & 17 & 10 & 55.1 & 41.2 & 25.7 \\
        11 & Kaio Jorge & Cruzeiro & 0.025 & [-0.18, 0.22] & 2,687 & 33 & 18 & 11 & 58.0 & 45.3 & 24.5 \\
        12 & Fabrício Bruno & Cruzeiro & 0.024 & [-0.17, 0.22] & 2,950 & 33 & 17 & 11 & 61.8 & 49.0 & 24.0 \\
        13 & Kaiki Bruno & Cruzeiro & 0.024 & [-0.17, 0.22] & 2,996 & 34 & 17 & 12 & 62.3 & 49.9 & 24.4 \\
        14 & Daniel & Mirassol & 0.024 & [-0.17, 0.22] & 3,265 & 37 & 18 & 12 & 65.5 & 60.8 & 23.9 \\
        15 & Varela & Flamengo & 0.024 & [-0.18, 0.22] & 1,326 & 20 & 13 & 5 & 36.1 & 31.6 & 25.6 \\
        16 & Pedro & Flamengo & 0.023 & [-0.17, 0.22] & 1,044 & 21 & 14 & 5 & 27.9 & 30.6 & 23.7 \\
        17 & Vitinho & Botafogo & 0.023 & [-0.17, 0.22] & 2,875 & 36 & 17 & 11 & 56.7 & 50.9 & 21.2 \\
        18 & Neto & Mirassol & 0.023 & [-0.18, 0.22] & 2,479 & 34 & 16 & 12 & 50.4 & 49.0 & 21.4 \\
        19 & Alex Sandro & Flamengo & 0.023 & [-0.17, 0.22] & 1,499 & 19 & 13 & 5 & 37.0 & 33.2 & 23.2 \\
        20 & Joaquin Piquerez & Palmeiras & 0.022 & [-0.17, 0.22] & 2,411 & 29 & 18 & 5 & 55.9 & 47.2 & 24.8 \\
        \bottomrule
    \end{tabular}
    \caption{
        \textbf{Top 20 players ranked by estimated skill under the P7 model for the 2025 season}.
        Skill estimates are posterior medians from 30,000 MCMC samples.
        Positive values indicate above-average skill relative to the league.
        Abbreviations: Min.\ = total minutes played, G = games, W = wins, D = draws, WS = weighted points (sum of team points weighted by player's share of match minutes), WG = weighted goals, WGD = weighted goal difference.
    }
    \label{tab:top20_2025}
\end{table}

The ranking is dominated by players from Flamengo (8 of top 20) and Cruzeiro (6), reflecting their strong performances as champion and third-place finisher, respectively.
The other players are also from clubs that finished in the top positions of the league, with Mirassol standing out with three players in their first season at \textit{Brasileirão Série A}: Jemmes, Daniel and Neto.
The top-ranked player, Agustin Rossi, has over 3,300 minutes with 23 wins in 37 matches, helping Flamengo to earn 78 points and the championship.
During his participation, Flamengo scored a weighted total of 75 goals with a goal difference of 51.

Just as the top rankings concentrate players from the best performing clubs, the bottom rankings reveal the other extreme, with players from relegated clubs.
Some of these players accumulated significant playing time on losing teams, resulting in skill estimates that reflect the team's poor performance.
For example, given Sport's relegation with only 17 points, their players unsurprisingly dominate the lowest positions, as shown in Table~\ref{tab:bottom20_2025}.

\begin{table}[htbp]
    \centering
    \tiny
    \begin{tabular}{rlcccrrrrrrrr}
        \toprule
        \# & Player & Team & Median & 95\% CI & Min. & G & W & D & WS & WG & WGD \\
        \midrule
        711 & Ewerthon & Juventude & -0.012 & [-0.21, 0.18] & 1,110 & 17 & 3 & 5 & 8.0 & 10.1 & -16.7 \\
        712 & Renato Kayzer & Vitória & -0.013 & [-0.21, 0.18] & 2,042 & 28 & 7 & 9 & 21.9 & 18.4 & -18.3 \\
        713 & Caique & Juventude & -0.014 & [-0.21, 0.18] & 2,054 & 28 & 7 & 6 & 18.8 & 19.5 & -20.1 \\
        714 & Ramon & Vitória & -0.014 & [-0.21, 0.18] & 1,949 & 23 & 7 & 5 & 25.1 & 15.1 & -20.1 \\
        715 & Chrystian & Sport & -0.015 & [-0.21, 0.18] & 1,546 & 33 & 2 & 11 & 6.8 & 10.7 & -18.7 \\
        716 & Ramon Menezes & Sport & -0.015 & [-0.21, 0.17] & 1,923 & 23 & 2 & 7 & 13.0 & 19.5 & -22.9 \\
        717 & Pablo & Sport & -0.015 & [-0.21, 0.18] & 1,098 & 20 & 1 & 5 & 2.3 & 10.1 & -21.9 \\
        718 & Alan Ruschel & Juventude & -0.015 & [-0.21, 0.18] & 1,275 & 20 & 3 & 4 & 7.1 & 8.7 & -18.5 \\
        719 & Caique França & Sport & -0.015 & [-0.21, 0.18] & 1,350 & 15 & 0 & 4 & 4.0 & 8.0 & -21.0 \\
        720 & Igor Carius & Sport & -0.015 & [-0.21, 0.18] & 1,822 & 26 & 2 & 8 & 10.2 & 13.6 & -21.5 \\
        721 & Luiz Gustavo & Juventude & -0.017 & [-0.21, 0.18] & 2,407 & 31 & 6 & 8 & 23.9 & 22.9 & -22.6 \\
        722 & Lucas & Sport & -0.018 & [-0.21, 0.17] & 904 & 11 & 0 & 0 & 0.0 & 6.1 & -23.9 \\
        723 & Luan Candido & Sport & -0.018 & [-0.21, 0.18] & 1,317 & 15 & 1 & 3 & 6.0 & 12.6 & -22.6 \\
        724 & Matheusinho & Sport & -0.019 & [-0.21, 0.18] & 1,702 & 23 & 2 & 6 & 10.2 & 16.3 & -23.8 \\
        725 & Rafael Thyere & Sport & -0.020 & [-0.21, 0.18] & 2,151 & 24 & 2 & 8 & 14.0 & 20.8 & -27.8 \\
        726 & Gabriel Vasconcelos & Sport & -0.021 & [-0.21, 0.17] & 2,160 & 24 & 2 & 7 & 13.0 & 20.0 & -28.0 \\
        727 & Jadson & Juventude & -0.022 & [-0.21, 0.17] & 2,920 & 35 & 9 & 7 & 29.9 & 30.6 & -31.4 \\
        728 & Matheus Alexandre & Sport & -0.023 & [-0.22, 0.17] & 1,782 & 25 & 2 & 7 & 9.0 & 16.1 & -28.1 \\
        729 & Lucas Lima & Sport & -0.028 & [-0.22, 0.16] & 2,549 & 32 & 1 & 9 & 10.8 & 20.6 & -36.5 \\
        730 & Christian Rivera & Sport & -0.028 & [-0.22, 0.16] & 2,818 & 33 & 2 & 8 & 14.0 & 24.3 & -39.5 \\
        \bottomrule
    \end{tabular}
    \caption{
        \textbf{Bottom 20 players ranked by estimated skill under the P7 model for the 2025 season}.
        Skill estimates are posterior medians from 30,000 MCMC samples.
        Abbreviations: Min.\ = total minutes played, G = games, W = wins, D = draws, WS = weighted points (sum of team points weighted by player's share of match minutes), WG = weighted goals, WGD = weighted goal difference.
    }
    \label{tab:bottom20_2025}
\end{table}

Beyond Sport, the bottom rankings include some players from Juventude, another club with a poor campaign that was relegated.
The lowest-ranked player, Christian Rivera, has over 2,800 minutes across 33 matches, during which Sport won only twice and earned 14 points, with a weighted goal difference of -39.5.
Some Sport players recorded zero wins across the entire season: Lucas (11 matches, weighted goal difference of -23.9) and Caique França (15 matches, weighted goal difference of -21.0).
In addition to Sport and Juventude players, the bottom 20 also include Ramon and Renato Kayzer, both from Vitória, a club that narrowly avoided relegation.

\subsection{2020--2025 Period}
\label{subsec:player_2020_2025}

While the single-season analysis provides a snapshot of player performance, extending the temporal window addresses two key limitations: small sample sizes for individual players and the confounding influence of short-term team dynamics.
The six-season period (2020--2025) increases both the sample size and the number of estimated player effects, with 1,975 unique players across 2,280 matches.
Before examining individual rankings, we first analyse the global parameter estimates.
Table~\ref{tab:player_global_2020_2025} presents the posterior summaries for this extended period.

\begin{table}[htbp]
    \centering
    \begin{tabular}{lcccc}
        \toprule
        Parameter & Median & Mean & Std & 95\% CI \\
        \midrule
        $\nu$ (home effect) & 0.324 & 0.324 & 0.027 & [0.270, 0.377] \\
        $\eta$ (common effect) & 0.016 & 0.016 & 0.021 & [$-0.025$, 0.057] \\
        $\alpha_{\text{synthetic}}$ & $-4.81$ & $-5.24$ & 2.42 & $(-\infty, -2.07]$ \\
        \bottomrule
    \end{tabular}
    \caption{
        \textbf{Posterior estimates for global parameters under the P7 model for the 2020--2025 period}.
        Estimates are computed from 30,000 MCMC samples.
    }
    \label{tab:player_global_2020_2025}
\end{table}

The home effect ($\nu = 0.324$) corresponds to a multiplicative factor of approximately 1.38 on expected goals, consistent with the range of club-level estimates in Table~\ref{tab:match_params} (log between 0.273 and 0.425, factor between 1.31 and 1.53).
The common effect parameter $\eta$ remains indistinguishable from zero, as in the single-season analysis.
The synthetic player parameter ($\alpha_{\text{synthetic}} = -4.81$) is considerably lower than the 2025 estimate ($-2.28$), and its credible interval now excludes zero by a wide margin.
In practical terms, this implies that receiving a red card is equivalent to replacing a player with one whose relative skill is $\exp(-4.81) \approx 0.008$ times the average.

With global parameters consistent with both the single-season and club-level analyses, we turn to individual player estimates.
The multi-season aggregation produces more stable skill coefficients, as players accumulate more minutes across multiple campaigns.
Tables~\ref{tab:top20_2020_2025} and~\ref{tab:bottom20_2020_2025} present the top and bottom 20 players ranked by posterior median log-skill, respectively.

\begin{table}[htbp]
    \centering
    \tiny
    \begin{tabular}{rlcccrrrrrrrr}
        \toprule
        \# & Player & Team & Median & 95\% CI & Min. & G & W & D & WS & WG & WGD \\
        \midrule
        1 & Weverton & Palmeiras & 0.132 & [-0.07, 0.33] & 16,448 & 183 & 105 & 44 & 359.0 & 308.0 & 153.5 \\
        2 & De Arrascaeta & Flamengo & 0.093 & [-0.10, 0.30] & 9,864 & 142 & 87 & 29 & 225.9 & 205.5 & 106.7 \\
        3 & Fabrício Bruno & Flamengo & 0.083 & [-0.12, 0.28] & 13,391 & 155 & 81 & 43 & 275.7 & 235.6 & 94.2 \\
        4 & Gustavo Gómez & Palmeiras & 0.083 & [-0.12, 0.29] & 13,786 & 157 & 86 & 39 & 290.7 & 236.2 & 105.5 \\
        5 & Murilo & Palmeiras & 0.082 & [-0.12, 0.28] & 8,930 & 102 & 60 & 27 & 202.5 & 167.3 & 96.4 \\
        6 & Joaquín Piquerez & Palmeiras & 0.079 & [-0.12, 0.28] & 9,453 & 112 & 68 & 24 & 216.5 & 179.2 & 99.9 \\
        7 & Marcos Rocha & Palmeiras & 0.076 & [-0.12, 0.27] & 9,723 & 133 & 75 & 34 & 213.0 & 178.8 & 91.1 \\
        8 & Raphael Veiga & Palmeiras & 0.076 & [-0.13, 0.28] & 11,203 & 164 & 91 & 37 & 236.4 & 200.3 & 93.0 \\
        9 & Léo Pereira & Flamengo & 0.067 & [-0.13, 0.27] & 11,332 & 136 & 72 & 35 & 231.1 & 207.6 & 84.0 \\
        10 & Zé Rafael & Palmeiras & 0.064 & [-0.14, 0.26] & 11,385 & 159 & 84 & 39 & 232.3 & 199.9 & 82.3 \\
        11 & Agustin Rossi & Flamengo & 0.063 & [-0.14, 0.26] & 7,830 & 87 & 49 & 22 & 169.0 & 149.0 & 74.0 \\
        12 & Éverson & Atlético-MG & 0.060 & [-0.14, 0.26] & 17,901 & 199 & 90 & 57 & 326.9 & 276.8 & 72.0 \\
        13 & Renê & Internacional & 0.058 & [-0.14, 0.26] & 9,691 & 128 & 67 & 33 & 197.1 & 152.3 & 61.0 \\
        14 & Rony & Palmeiras & 0.056 & [-0.14, 0.25] & 11,222 & 181 & 90 & 49 & 221.0 & 190.4 & 73.1 \\
        15 & Léo Ortiz & RB Bragantino & 0.055 & [-0.15, 0.25] & 12,565 & 149 & 65 & 48 & 226.6 & 202.7 & 64.7 \\
        16 & Pedro & Flamengo & 0.052 & [-0.15, 0.25] & 9,030 & 157 & 87 & 33 & 184.7 & 170.2 & 65.9 \\
        17 & Bruno Henrique & Flamengo & 0.052 & [-0.15, 0.25] & 10,385 & 148 & 81 & 35 & 211.8 & 182.4 & 64.7 \\
        18 & Ayrton Lucas & Flamengo & 0.052 & [-0.15, 0.25] & 8,812 & 116 & 60 & 26 & 174.8 & 159.9 & 61.6 \\
        19 & Gerson & Flamengo & 0.050 & [-0.15, 0.25] & 8,946 & 109 & 61 & 25 & 192.2 & 173.4 & 63.5 \\
        20 & Victor Cuesta & Internacional & 0.050 & [-0.15, 0.25] & 12,449 & 141 & 63 & 40 & 226.6 & 192.8 & 55.3 \\
        \bottomrule
    \end{tabular}
    \caption{
        \textbf{Top 20 players ranked by estimated skill under the P7 model for the 2020--2025 period}.
        Skill estimates are posterior medians from 30,000 MCMC samples.
        Abbreviations: Min.\ = total minutes played, G = games, W = wins, D = draws, WS = weighted points (sum of team points weighted by player's share of match minutes), WG = weighted goals, WGD = weighted goal difference.
    }
    \label{tab:top20_2020_2025}
\end{table}

\begin{table}[htbp]
    \centering
    \tiny
    \begin{tabular}{rlcccrrrrrrrr}
        \toprule
        \# & Player & Team & Median & 95\% CI & Min. & G & W & D & WS & WG & WGD \\
        \midrule
        1,956 & Victor Luis & Botafogo & -0.026 & [-0.22, 0.17] & 6,632 & 104 & 23 & 31 & 67.1 & 74.1 & -32.7 \\
        1,957 & Lucas & América-MG & -0.026 & [-0.22, 0.17] & 7,119 & 94 & 27 & 20 & 89.6 & 78.8 & -34.7 \\
        1,958 & Rodrigo Alves & Juventude & -0.026 & [-0.22, 0.16] & 2,788 & 35 & 3 & 11 & 17.9 & 23.0 & -34.4 \\
        1,959 & Pedro Raul & Ceará & -0.026 & [-0.22, 0.16] & 8,898 & 123 & 27 & 39 & 101.1 & 93.5 & -37.9 \\
        1,960 & Victor Hugo & Bahia & -0.027 & [-0.22, 0.16] & 10,710 & 124 & 37 & 27 & 132.6 & 128.9 & -36.8 \\
        1,961 & Shaylon & Atlético-GO & -0.028 & [-0.22, 0.17] & 6,143 & 103 & 26 & 28 & 65.1 & 63.3 & -36.6 \\
        1,962 & Rafael Thyere & Sport & -0.028 & [-0.22, 0.17] & 5,567 & 66 & 14 & 20 & 56.0 & 49.1 & -37.2 \\
        1,963 & Lucas Halter & Goiás & -0.029 & [-0.22, 0.16] & 8,664 & 104 & 33 & 27 & 112.1 & 94.5 & -39.1 \\
        1,964 & Christian Rivera & Sport & -0.029 & [-0.22, 0.16] & 2,818 & 33 & 2 & 8 & 14.0 & 24.3 & -39.5 \\
        1,965 & Gilberto & Bahia & -0.030 & [-0.22, 0.16] & 8,156 & 133 & 37 & 30 & 93.8 & 99.0 & -38.8 \\
        1,966 & Matheus & Cuiabá & -0.030 & [-0.22, 0.16] & 8,629 & 111 & 27 & 26 & 96.8 & 91.8 & -37.9 \\
        1,967 & Gabriel Baralhas & Atlético-GO & -0.031 & [-0.23, 0.16] & 9,088 & 129 & 37 & 35 & 110.1 & 96.7 & -44.6 \\
        1,968 & Éder & Atlético-GO & -0.034 & [-0.22, 0.16] & 10,617 & 126 & 37 & 35 & 136.3 & 111.3 & -46.0 \\
        1,969 & Tadeu & Goiás & -0.035 & [-0.22, 0.15] & 9,450 & 105 & 26 & 31 & 109.0 & 107.0 & -48.0 \\
        1,970 & Matheusinho & Sport & -0.036 & [-0.23, 0.16] & 4,960 & 85 & 20 & 20 & 52.5 & 55.4 & -41.3 \\
        1,971 & Isidro Pitta & Cuiabá & -0.037 & [-0.23, 0.15] & 7,161 & 118 & 31 & 27 & 66.6 & 64.7 & -48.8 \\
        1,972 & Kevin & Avaí & -0.037 & [-0.23, 0.15] & 5,213 & 65 & 11 & 15 & 45.4 & 51.6 & -47.6 \\
        1,973 & Lucas Lima & Sport & -0.044 & [-0.24, 0.15] & 8,823 & 137 & 32 & 40 & 91.4 & 93.5 & -54.0 \\
        1,974 & Jadson & Juventude & -0.050 & [-0.24, 0.14] & 10,300 & 131 & 30 & 38 & 113.0 & 114.7 & -66.4 \\
        1,975 & Gabriel Vasconcelos & Coritiba & -0.051 & [-0.24, 0.14] & 9,110 & 102 & 23 & 25 & 93.5 & 104.2 & -71.4 \\
        \bottomrule
    \end{tabular}
    \caption{
        \textbf{Bottom 20 players ranked by estimated skill under the P7 model for the 2020--2025 period}.
        Skill estimates are posterior medians from 30,000 MCMC samples.
        Abbreviations: Min.\ = total minutes played, G = games, W = wins, D = draws, WS = weighted points (sum of team points weighted by player's share of match minutes), WG = weighted goals, WGD = weighted goal difference.
    }
    \label{tab:bottom20_2020_2025}
\end{table}

Two points immediately stand out from the rankings.
First, all credible intervals (both top and bottom and in all periods) include zero, meaning the model cannot statistically distinguish any individual player from the league-wide average.
Second, the top 20 is heavily concentrated in two clubs: Palmeiras and Flamengo, with 8 players each.

Despite the statistical indistinguishability, examining the ranking by posterior median reveals interpretable patterns.
In this period, Palmeiras was champion twice, in 2022 and 2023, and finished in the top half of the table in the other seasons: 7th in 2020, 3rd in 2021 and 2nd in both 2024 and 2025.
Flamengo also won the championship twice, in 2020 and 2025, with good results in the other seasons: 2nd in 2021, 5th in 2022, 4th in 2023 and 3rd in 2024.

The four remaining players in the top 20 (Éverson, Renê, Léo Ortiz, and Victor Cuesta) also accumulated significant minutes in successful environments, though across multiple clubs.
Éverson won the 2021 title with Atlético-MG and finished in the top three twice more.
Renê played for Flamengo (2020--2021), Internacional (2022--2024), and Fluminense (2025), consistently participating in top-half finishes.
Léo Ortiz spent four seasons at RB Bragantino before transferring to Flamengo in 2024, while Victor Cuesta played for Internacional, Botafogo, and Bahia, all clubs that achieved mid-to-upper table finishes during his tenure.
Conversely, the bottom 20 includes players from a broader set of teams that experienced relegation or persistent struggles: Coritiba, Juventude, Sport, Cuiabá, Avaí, Goiás, and Atlético Goianiense.
These patterns indicate that posterior rankings reflect sustained exposure to winning or losing contexts rather than isolated individual excellence or deficiency.

This interpretation is reinforced by the near-perfect correlation ($r = 0.99$) between posterior medians and weighted goal difference (WGD).
Since WGD aggregates team-level outcomes (goals scored minus goals conceded) weighted by each player's share of match minutes, this correlation reveals that estimated skill largely captures whether a player recorded minutes in winning or losing environments.
In the end, the model effectively transforms collective team performance into individual-level estimates, with limited capacity to isolate marginal contributions.

This structural dependence has notable consequences for certain player profiles.
Goalkeepers, who typically play full matches whenever fielded, are particularly susceptible to extreme rankings because their WGD closely tracks their team's aggregate performance.
Among the top 20, three goalkeepers appear: Weverton (1st), Agustin Rossi (11th), and Éverson (12th).
Among the bottom 20, Tadeu (1,969th) and Gabriel Vasconcelos (1,975th) occupy similarly extreme positions.
The pattern extends beyond goalkeepers to any player who consistently appears in the starting lineup of a club and plays a considerable fraction of the match.

The case of Pedro Raul illustrates this dynamic.
Despite being the runner-up top scorer of the 2022 \textit{Brasileirão} while playing for Goiás (13th place in that season) and scoring 11 goals for Ceará in 2025, he ranks 1,959th out of 1,975 players in the six-season analysis.
Across 123 matches and nearly 9,000 minutes, he participated in only 27 wins and recorded a weighted goal difference of $-37.9$.
These numbers reflect the poor performance of teams like Goiás and Ceará rather than his individual ability.
Since the model observes only match outcomes, his individual attacking output cannot offset the collective struggles of his teams.

These findings echo the club-level analysis of Section~\ref{sec:club_applications}.
Just as team strength estimates reflect aggregate match outcomes, player-level estimates reflect the teams in which players played.
The decomposition from clubs to players increases model dimensionality but does not resolve the fundamental challenge of separating individual ability from collective performance.
Future work incorporating player-specific actions or event data may offer a path toward more granular attribution.
